% ============================================================
% 5. PRUEBAS
% ============================================================
\section{Pruebas}
\label{sec:pruebas}

\subsection{Diseño de las Pruebas}
\label{sec:diseno-pruebas}

Cada capacidad de la solución se verifica con una prueba que sigue la
misma estructura: se comprueba y se registra el \emph{estado anterior}
de la máquina con órdenes ejecutadas fuera de la consola ---terminal de
la VM, de otra VM o del anfitrión---, se realiza la operación desde la
interfaz web y se vuelve a comprobar el \emph{estado posterior} con las
mismas órdenes. La comprobación externa es deliberada: si el estado
posterior se leyera solo en la interfaz, una respuesta incorrecta del
servidor podría pasar por un éxito. Cada prueba se graba en video desde
el navegador y la terminal al mismo tiempo, de modo que el estado
anterior y el posterior quedan visibles en el mismo registro.

El entorno es el de la \cref{fig:escenario}: el anfitrión con
VirtualBox~7.2 y la consola; el servidor S1 (\texttt{192.168.1.29}),
preparado desde el disco de multiconexión con Debian~13.7.0 en modo
puente y el servicio de imágenes en el puerto~8000, que se usa además
para las pruebas de preparación del sistema; el servidor S2
(\texttt{192.168.1.30}), con la misma distribución y el mismo servicio,
que completa el back-end y sirve para el cambio de máquina; y el
balanceador B (\texttt{192.168.1.32}), con HAProxy sobre Debian~13.7.0.
En el anfitrión se usa \texttt{oha}~1.16 para generar carga sobre el
balanceador. Las direcciones IP se obtienen por DHCP y no se fijan en la
prueba: la consola las descubre en tiempo de ejecución.

Las \cref{tab:diseno-pruebas,tab:diseno-pruebas-b}
resumen las veintidós pruebas diseñadas. Las pruebas P05 y P16 incluyen, además
del camino correcto, el camino de error: un paquete que no existe en los
repositorios y la baja repetida de un servidor del back-end; P22 añade la
carga concurrente con \texttt{oha} para observar el reparto entre los dos
servidores.

\begin{table*}[!tp]
  \centering
  \caption{Diseño de las pruebas (parte 1 de 2): estado anterior,
  operación y estado posterior esperado de cada elemento de la solución.}
  \label{tab:diseno-pruebas}
  \scriptsize
  \begin{tabularx}{\textwidth}{@{}l >{\raggedright\arraybackslash}p{2.3cm}
      >{\raggedright\arraybackslash}X >{\raggedright\arraybackslash}X
      >{\raggedright\arraybackslash}X
      >{\raggedright\arraybackslash}p{1.0cm}@{}}
    \hline
    ID & Capacidad & Estado anterior & Operación desde la consola &
    Estado posterior esperado & Video \\
    \hline
    P01 & Registro de la máquina & S2 apagada y sin registrar; sin
    llave de root en el anfitrión. & Encender desde el selector y
    registrar con el usuario y la contraseña de la VM. &
    \texttt{/api/session} informa la IP y el conector \texttt{ssh}; la
    máquina aparece registrada y no vuelve a pedir credenciales. & V01 \\
    P02 & Diagnóstico del registro & VM con Guest Additions~6.0 sobre
    VirtualBox~7.2, cuyo \texttt{guestcontrol} rechaza la contraseña. &
    Enviar las credenciales en el modal y abrir el diagnóstico. &
    Mensaje comprensible y diagnóstico con versión de Guest Additions,
    puerto 22 y usuario con sesión; el registro se completa por SSH con
    \texttt{sshpass}. & V02 \\
    P03 & Instalar paquete & \texttt{dpkg-query -W htop} responde que no
    está instalado. & Formulario de paquetes: instalar \texttt{htop}. &
    \texttt{htop} figura como \texttt{install ok installed}, aparece en
    el listado y la ficha muestra el código de salida cero. & V03 \\
    P04 & Eliminar paquete & \texttt{htop} instalado. & Eliminar
    \texttt{htop} y confirmar en el aviso. & \texttt{htop} ya no está
    instalado, desaparece del listado y la ficha describe la orden
    ejecutada. & V04 \\
    P05 & Error de APT & Paquete \texttt{paquete-inexistente} no
    disponible en los repositorios. & Intentar instalarlo. & Respuesta
    \texttt{500} con mensaje comprensible, código de salida de APT y
    \texttt{stderr} en el detalle técnico; la máquina queda sin cambios.
    & V05 \\
    P06 & Llave de \texttt{root} & Sin
    \texttt{/etc/ssh/\allowbreak sshd\_config.d/\allowbreak
    99-vmcons.conf} ni llave en el
    anfitrión. & Generar la llave y descargar la privada. &
    \texttt{authorized\_keys} con permisos \texttt{700}/\texttt{600};
    \texttt{sshd -T} informa \texttt{passwordauthentication no} y
    \texttt{permitrootlogin prohibit-password}; la descarga entrega un
    archivo adjunto y la interfaz recomienda resguardarlo. & V06 \\
    P07 & Acceso por contraseña & Contraseña de \texttt{root}
    deshabilitada por P06. & \texttt{ssh -o \allowbreak
    PubkeyAuthentication=no \allowbreak root@IP}. & El intento por
    contraseña se rechaza y el acceso con \texttt{-i llave} funciona.
    & V07 \\
    P08 & \texttt{/etc/issue} & Banner original de Debian, que la
    consola muestra como texto ajeno con las casillas desactivadas. &
    Escribir un mensaje y marcar nombre de host, fecha, hora e IP. &
    \texttt{cat /etc/issue} muestra las etiquetas y las secuencias de
    \texttt{agetty}; la consola de inicio de sesión las expande. & V08 \\
    P09 & Nombre de host y \texttt{/etc/hosts} & \texttt{hostname}
    devuelve \texttt{debian} y \texttt{/etc/hosts} no menciona la IP de
    la VM. & Asignar \texttt{vmcons-lab} con la casilla de actualizar
    \texttt{/etc/hosts}. & \texttt{hostname} devuelve
    \texttt{vmcons-lab}, \texttt{/etc/hosts} conserva
    \texttt{localhost} y agrega «IP nombre»; el cambio persiste tras
    reiniciar. & V09 \\
    P10 & Crear usuario con \texttt{sudo} y llaves & \texttt{id alumno}
    responde que no existe la cuenta. & Crear \texttt{alumno} con las
    casillas de \texttt{sudo} y de llaves, y descargar la privada. &
    UID entre 1000 y 65533, grupo \texttt{sudo},
    \texttt{authorized\_keys} del usuario con \texttt{700}/\texttt{600}
    y contraseña bloqueada; la cuenta aparece en el listado. & V10 \\
    P11 & Eliminar usuario & \texttt{alumno} existe y tiene
    \texttt{/home/alumno}. & Eliminar con la casilla de borrar archivos.
    & \texttt{id alumno} falla, \texttt{/home/alumno} ya no existe y la
    llave del anfitrión se elimina; la cuenta sale del listado. & V11 \\
    \hline
  \end{tabularx}
\end{table*}

\begin{table*}[!tp]
  \centering
  \caption{Diseño de las pruebas (parte 2 de 2): estado anterior,
  operación y estado posterior esperado de cada elemento de la solución.}
  \label{tab:diseno-pruebas-b}
  \scriptsize
  \begin{tabularx}{\textwidth}{@{}l >{\raggedright\arraybackslash}p{2.3cm}
      >{\raggedright\arraybackslash}X >{\raggedright\arraybackslash}X
      >{\raggedright\arraybackslash}X
      >{\raggedright\arraybackslash}p{1.0cm}@{}}
    \hline
    ID & Capacidad & Estado anterior & Operación desde la consola &
    Estado posterior esperado & Video \\
    \hline
    P12 & Desplegar una unidad & La aplicación del catálogo no está en
    la VM ni tiene unidad. & Subir el ejecutable y el archivo
    \texttt{.service} con modo \texttt{0755} y propietario. &
    \texttt{systemctl cat} muestra la unidad, \texttt{daemon-reload}
    queda aplicado y el listado la informa inactiva y deshabilitada.
    & V12 \\
    P13 & Ciclo de vida de la unidad & Unidad inactiva y deshabilitada.
    & \texttt{start}, \texttt{enable}, \texttt{restart},
    \texttt{stop} y \texttt{disable} desde los botones. &
    \texttt{systemctl is-active} y \texttt{is-enabled} reflejan cada
    paso; \texttt{status} muestra un PID distinto tras el reinicio.
    & V13 \\
    P14 & Script de la precondición e) & Unidad del script inactiva y
    archivo de registro sin líneas del día. & Iniciar y reiniciar la
    unidad. & Cada ejecución agrega una línea con la fecha y la hora;
    tras el reinicio hay dos líneas nuevas y un PID distinto. & V14 \\
    P15 & Agregar servidores al back-end & \texttt{backend}
    \texttt{vmcons\_app} vacío en el balanceador. & Agregar
    \texttt{192.168.1.29:8000} y \texttt{192.168.1.30:8000}. & Dos
    líneas \texttt{server srv\_\dots\ check}, configuración válida con
    \texttt{haproxy -c}, servicio recargado y \texttt{curl} al
    balanceador atendido por los dos servidores. & V15 \\
    P16 & Eliminar servidores del back-end & S1 y S2 en el back-end
    del balanceador. & Eliminar S1 y después S2. & La línea de S1
    desaparece y el tráfico solo llega a S2; sin servidores el
    balanceador responde \texttt{503} y repetir la baja da
    \texttt{404}. & V16 \\
    P17 & Historial y persistencia & Base con operaciones anteriores de
    la sesión de pruebas. & Ejecutar tres operaciones y consultar la
    pestaña Historial con filtros de texto, fecha, máquina y módulo;
    exportar a CSV. & Las tres operaciones aparecen con descripción
    legible; los filtros acotan el resultado; el CSV coincide y las
    filas sobreviven al reinicio de la consola. & V17 \\
    P18 & Cambio de máquina vinculada & S1 vinculada y listados
    cargados en varias pestañas. & Vincular S2 desde el selector.
    & La barra muestra la IP de S2 y los listados se recargan sin
    conservar filas de S1. & V18 \\
    P19 & Estado de HAProxy & HAProxy con la configuración por omisión,
    sin bloques \texttt{backend}. & Abrir el módulo HAProxy y consultar
    el estado. & La tarjeta indica que requiere configuración con el
    problema «no tiene bloques backend»; el estado informa el binario,
    la configuración válida, el servicio y el puerto. & V19 \\
    P20 & Configuración rápida & HAProxy sin \texttt{frontend} ni
    \texttt{backend} propios. & Pulsar \textbf{Configuración rápida}.
    & \texttt{haproxy.cfg} con \texttt{frontend vmcons\_front} en el
    puerto~80 y \texttt{backend vmcons\_app}; copia
    \texttt{.bak-\dots}, \texttt{haproxy -c} válida, servicio activo y
    tarjeta con la URL de entrada. & V20 \\
    P21 & Métricas de los servidores & S1 y S2 agregados al back-end.
    & Abrir el detalle de cada servidor y recargar. & Estado
    \texttt{UP}, peso, activo, sesiones, cola, bytes enviados y
    recibidos, \texttt{L4OK}, duración del chequeo y tiempos; los
    totales crecen con el tráfico. & V21 \\
    P22 & Balanceo bajo carga y URL de entrada & Totales similares en S1
    y S2. & Copiar la URL del balanceador y ejecutar
    \texttt{oha --no-tui -n 600 -c 30 http://192.168.1.32/}. & Todas las
    respuestas son \texttt{2xx}; ambos servidores suman unas 300
    peticiones y el máximo de sesiones sube; la URL copiada abre la
    aplicación. & V22 \\
    \hline
  \end{tabularx}
\end{table*}

\subsection{Resultados y Evidencias}
\label{sec:resultados-pruebas}

\subsubsection{Verificación Automática}
\label{sec:resultados-automaticos}

Las pruebas automatizadas verifican la lógica que no depende de una
máquina real: la validación de entradas, la traducción de errores, el
descubrimiento y la selección del conector, la construcción de las
órdenes, el manejo de los recursos HTTP, la persistencia, la caché y el
historial. La \cref{tab:cobertura} resume el resultado; las órdenes se
ejecutan con las tareas del proyecto (\texttt{task test},
\texttt{task lint}, \texttt{task openapi} y \texttt{pnpm test}).

\begin{table}[!t]
  \centering
  \caption{Resultado de las pruebas automatizadas.}
  \label{tab:cobertura}
  \footnotesize
  \begin{tabularx}{\columnwidth}{@{}l >{\raggedleft\arraybackslash}X@{}}
    \hline
    Paquete & Cobertura de sentencias \\
    \hline
    \texttt{core/domain} & 100.0\,\% \\
    \texttt{adapters/localproc} & 100.0\,\% \\
    \texttt{httpcache} & 100.0\,\% \\
    \texttt{handlers} & 98.5\,\% \\
    \texttt{services} & 94.9\,\% \\
    \texttt{adapters/ssh} & 100.0\,\% \\
    \texttt{adapters/vbox} & 100.0\,\% \\
    Ejecutable (\texttt{main}) & 82.5\,\% \\
    \hline
  \end{tabularx}
\end{table}

Las pruebas automatizadas terminan sin fallos y el análisis estático no
reporta incidencias. Los paquetes internos alcanzan el 100\,\% salvo
\texttt{handlers} (98.5\,\%) y \texttt{services} (94.9\,\%), donde las
sentencias restantes son ramas de error de restauración, cifrado y cierre
de sesión que exigen fallos encadenados. Los adaptadores
y los servicios se prueban con ejecutores y sesiones falsos que devuelven
salidas realistas de \texttt{VBoxManage} y \texttt{ssh}, y los
manejadores con peticiones HTTP simuladas; las pruebas de campo P01 a P22
completan lo que solo tiene sentido frente a un hipervisor real. El
paquete del ejecutable queda en el 82.5\,\%: su función de arranque y el
servidor que queda escuchando no se pueden ejercitar en una prueba
unitaria sin abrir un puerto. El contrato OpenAPI se valida sin errores y
ninguna de sus rutas devuelve \texttt{501}.

\subsubsection{Pruebas de Campo}
\label{sec:resultados-campo}

Las veintidós pruebas de las \cref{tab:diseno-pruebas,tab:diseno-pruebas-b} se ejecutan sobre
el escenario y se graban en video, con el nombre \texttt{V01} a
\texttt{V22} según la tabla. Cada grabación empieza con la comprobación
del estado anterior en la terminal, continúa con la operación en la
interfaz ---el formulario, la validación y la ficha del panel de
resultado--- y termina con la comprobación del estado posterior en la
misma terminal, de modo que la evidencia cubre el ciclo completo y no
solo la respuesta de la consola.

La ejecución del 26 de octubre de 2026 cubrió las veintidós pruebas
sobre los servidores S1 (\texttt{192.168.1.29}) y S2
(\texttt{192.168.1.30}) y el balanceador B (\texttt{192.168.1.32}). Las
máquinas se renombraron a S1, S2 y B para la grabación y el DHCP les
asignó esas direcciones; el registro de la ejecución conserva los nombres
y las direcciones anteriores. Las pruebas P01 a P18 se ejecutaron con el
arnés automatizado, que comprobó
el estado anterior por SSH, ejecutó la operación y verificó el estado
posterior con la misma conexión; P19 a P22 se verificaron en la misma
sesión desde la interfaz y el socket de estadísticas, con \texttt{oha}
como generador de carga. Las veintidós terminaron conformes; la
\cref{tab:resultados-campo} resume lo observado
y el detalle completo, con las salidas y los códigos de cada llamada,
está en \texttt{docs/resultados-pruebas.md}. Las pruebas restauran el
estado que modifican: la llave privada del anfitrión, el nombre de host,
\texttt{/etc/issue}, la unidad temporal y los servidores del
balanceador quedaron como estaban antes de la ejecución, y la
configuración de HAProxy se recuperó de la copia guardada antes de P19.

\begin{table*}[!t]
  \centering
  \caption{Resultados observados en las pruebas de campo.}
  \label{tab:resultados-campo}
  \scriptsize
  \begin{tabularx}{\textwidth}{@{}l >{\hsize=0.8\hsize}X >{\hsize=0.9\hsize}X >{\hsize=1.3\hsize}X c@{}}
    \hline
    Prueba & Antes & Operación & Después & Resultado \\
    \hline
    P01 & S2 sin registrar & \texttt{POST /api/session/register} & \texttt{registrada=true}; Guest Additions 6.0.0 & Conforme \\
    P02 & S2 registrada & \texttt{GET /api/session/diagnostico} & Red puente, IP \texttt{.30}, SSH abierto y aviso de Guest Additions & Conforme \\
    P03 & \texttt{tree} no instalado & \texttt{POST /api/packages} & \texttt{install ok installed} & Conforme \\
    P04 & \texttt{tree} instalado & \texttt{DELETE /api/packages/tree} & \texttt{dpkg-query}: no encontrado & Conforme \\
    P05 & \texttt{tree} eliminado & Alta de paquete inexistente & \texttt{500} con código 1 y \texttt{stderr} de APT & Conforme \\
    P06 & Sin llave privada en el anfitrión & \texttt{POST} y \texttt{GET /api/keys/root} & SSH con la llave nueva; \texttt{passwordauthentication no} & Conforme \\
    P07 & Contraseña deshabilitada & SSH sin autenticación por llave & \texttt{Permission denied (publickey)} & Conforme \\
    P08 & \texttt{/etc/issue} original & \texttt{PUT /api/issue} & Mensaje con escapes \texttt{\textbackslash n}, \texttt{\textbackslash d}, \texttt{\textbackslash t} y \texttt{\textbackslash 4} & Conforme \\
    P09 & Nombre de host \texttt{base} & \texttt{PUT /api/hostname} & \texttt{vmcons-lab} y entrada en \texttt{/etc/hosts} & Conforme \\
    P10 & Usuario \texttt{alumno} inexistente & \texttt{POST /api/users} & UID 1001, grupo \texttt{sudo}, llave 0600 y descarga & Conforme \\
    P11 & \texttt{alumno} con directorio & \texttt{DELETE /api/\allowbreak users/\allowbreak alumno?files=true} & Usuario y \texttt{/home/alumno} eliminados & Conforme \\
    P12 & Unidad inexistente & \texttt{PUT /api/units/vmcons-demo.service} & Unidad creada; \texttt{daemon-reload} & Conforme \\
    P13 & Unidad inactiva & Activar, habilitar, reiniciar y detener & \texttt{active}, \texttt{enabled}, \texttt{202}, \texttt{inactive} & Conforme \\
    P14 & Sin ejecuciones del script & Iniciar y reiniciar la unidad & Cinco líneas con fecha y hora ISO-8601 & Conforme \\
    P15 & Back-end sin servidores & \texttt{POST /api/\allowbreak haproxy/\allowbreak backend/\allowbreak servers} con S1 y S2 & Líneas \texttt{srv\_192\_168\_1\_29\_8000} y \texttt{srv\_192\_168\_1\_30\_8000}; \texttt{haproxy -c} válida y \texttt{curl} 200 & Conforme \\
    P16 & S1 y S2 en el back-end & \texttt{DELETE /api/\allowbreak haproxy/\allowbreak backend/\allowbreak servers} & S1 fuera y solo S2 atiende; sin servidores \texttt{503} y repetir 404 & Conforme \\
    P17 & Operaciones anteriores & \texttt{GET /api/historial?q=paquetes} & Operaciones con descripción legible & Conforme \\
    P18 & S1 vinculada & \texttt{PUT /api/session} con S2 & IP \texttt{.29} sustituida por \texttt{.30} & Conforme \\
    P19 & Configuración sin \texttt{backend} & \texttt{GET /api/haproxy/status} & Problema «no tiene bloques backend»; binario, configuración, servicio y puerto informados & Conforme \\
    P20 & Sin \texttt{frontend} ni \texttt{backend} & \texttt{POST /api/haproxy/configurar} & \texttt{vmcons\_front} en el puerto 80, \texttt{vmcons\_app}, copia \texttt{.bak-\dots} y servicio activo & Conforme \\
    P21 & S1 y S2 agregados & \texttt{GET /api/\allowbreak haproxy/\allowbreak backend/\allowbreak servers} & \texttt{UP}, peso 1, \texttt{L4OK}, sesiones, bytes y más de 13\,000 peticiones acumuladas & Conforme \\
    P22 & Totales similares & \texttt{oha --no-tui -n 600 -c 30} & 100\,\% de respuestas \texttt{2xx}; unos 300 registros más en cada servidor & Conforme \\
    \hline
  \end{tabularx}
\end{table*}

Los videos se entregan junto con el trabajo, fuera de este documento,
con el nombre que indica la tabla. Las capturas de la interfaz que
acompañan al informe están en la Sección~IV, donde se explica el módulo
correspondiente.

\begin{table}[!t]
  \centering
  \caption{Videos de las pruebas de campo, entregados por separado.}
  \label{tab:evidencias}
  \footnotesize
  \begin{tabularx}{\columnwidth}{@{}l X@{}}
    \hline
    Video & Contenido \\
    \hline
    V01--V02 & Registro y diagnóstico de la máquina \\
    V03--V05 & Paquetes: instalar, eliminar y error de APT \\
    V06--V07 & Llave de \texttt{root} y bloqueo de la contraseña \\
    V08--V09 & \texttt{/etc/issue} y nombre de host \\
    V10--V11 & Creación y eliminación de usuarios \\
    V12--V14 & Despliegue y ciclo de vida con systemd \\
    V15--V16 & Back-end de HAProxy: alta y baja \\
    V17--V18 & Historial y cambio de máquina \\
    V19--V22 & HAProxy: estado, configuración rápida, métricas y carga \\
    \hline
  \end{tabularx}
\end{table}

% Las tablas de esta sección no deben aparecer dentro de las conclusiones.
\FloatBarrier
